\thispagestyle{plain}
\begin{center}
    \Large \textbf{\uppercase{Abstract}}
\end{center}

\vspace{3\baselineskip}

\noindent
Person re-identification (Re-ID) aims to retrieve images of a specific individual across non-overlapping camera views. While supervised methods achieve remarkable performance, they require expensive identity-annotated datasets. Unsupervised methods address this by learning discriminative representations without identity labels, typically through iterative clustering and pseudo label-based training.

This thesis presents a reproduction study of the \textbf{Person Intrinsic Semantic Learning (PISL)} framework, proposed by Tao~\textit{et~al.} in IEEE Transactions on Image Processing (2024). PISL introduces three key innovations: (1)~a Spatial Diffusion Model (SDM) that generates semantically meaningful body/vehicle patches through a denoising diffusion process over spatial transformer parameters, (2)~a Semantic Controlled Diffusion (SCD) loss that guides the denoising direction, and (3)~a Patch Semantic Consistency (PSC) loss that refines pseudo labels by ensuring consistency between local and global features. The framework additionally incorporates camera-aware contrastive learning to mitigate cross-camera domain gaps.

We reproduce the camera-aware PISL model on two benchmark datasets --- Market-1501 (person Re-ID) and VeRi-776 (vehicle Re-ID) --- using six NVIDIA RTX~2080~Ti GPUs. On VeRi-776, our reproduction achieves \textbf{41.7\% mAP} and \textbf{87.1\% Rank-1} accuracy (42.9\% mAP with re-ranking), closely matching the paper's reported 43.1\% mAP. On Market-1501, we achieve \textbf{77.7\% mAP} and \textbf{91.0\% Rank-1} accuracy (88.6\% mAP with re-ranking). Comprehensive qualitative analysis through attention heatmaps, retrieval ranking visualizations, and t-SNE feature embeddings validates the framework's ability to learn semantically structured representations.

\vspace{\baselineskip}

\noindent
\textbf{Keywords}: Unsupervised person re-identification, diffusion model, spatial transformer network, pseudo label refinement, contrastive learning, semantic consistency.
